\documentclass[10pt,a4paper]{article}
\usepackage[margin=2.0cm,top=2.0cm,bottom=2.2cm]{geometry}
\usepackage{fontspec}
\usepackage{amsmath,amssymb}
\usepackage{booktabs}
\usepackage{tabularx}
\usepackage{fancyvrb}
\usepackage{xcolor}
\usepackage{enumitem}
\usepackage{fancyhdr}
\usepackage[colorlinks=true,linkcolor=blue!45!black,urlcolor=blue!45!black]{hyperref}

\setmainfont{DejaVu Sans}[Scale=0.88]
\setsansfont{DejaVu Sans}[Scale=0.88]
\setmonofont{DejaVu Sans Mono}[Scale=0.82]

\definecolor{rulegray}{gray}{0.72}
\newenvironment{quoting}%
  {\begin{list}{}{\setlength{\leftmargin}{1.4em}\setlength{\rightmargin}{1.4em}%
   \setlength{\topsep}{5pt}}\item[]\itshape}%
  {\end{list}}
\setlength{\parindent}{0pt}
\setlength{\parskip}{0.55em}
\renewcommand{\arraystretch}{1.18}

\pagestyle{fancy}\fancyhf{}
\fancyhead[L]{\footnotesize\sffamily Parameterised water balance --- design for critique}
\fancyhead[R]{\footnotesize\sffamily 2026-08-24}
\fancyfoot[C]{\footnotesize\thepage}
\renewcommand{\headrulewidth}{0.3pt}

\title{\bfseries\sffamily A parameterised water balance\\driven by a frozen foundation model}
\author{\sffamily Soil moisture project --- Session 30}
\date{\sffamily 2026-08-24}

\begin{document}
\maketitle
\thispagestyle{fancy}
{\small\tableofcontents}
\clearpage

\textbf{Design document for critique. Nothing built. No GPU spent.}

Session 30, 2026-08-24. Supersedes nothing yet — this is an alternative to
§31/§32's per-location processor, not a replacement for it. §32's terrain
pipeline is a dependency of both and is unaffected.

\medskip\hrule\medskip

\subsection*{The idea in one sentence}
\addcontentsline{toc}{subsection}{The idea in one sentence}

\begin{quoting}
The network does not predict soil moisture. It predicts \textbf{the properties of the ground}. A differentiable water balance predicts soil moisture from those properties and the weather.
\end{quoting}

\medskip\hrule\medskip

\section*{0. FEASIBILITY}
\addcontentsline{toc}{section}{0. FEASIBILITY}

This section is first because it is the part that can kill the design, and
because two of its findings changed the architecture rather than merely
assessing it.

\subsection*{0.1 What the data actually gives}
\addcontentsline{toc}{subsection}{0.1 What the data actually gives}

Measured 2026-08-24 from \texttt{csvs/station\_splits.csv}:

\begin{center}
\small
\begin{tabularx}{\linewidth}{lX}
\toprule
\textbf{quantity} & \textbf{value} \\
\midrule
stations & 993 — train 683, val 91, oos 219 \\
station-years & 7,164 (mean 7.21, median 6.90, range 3–12) \\
station-days & 2,614,787 \\
location groups & 906 \\
Köppen classes / networks & 25 / 33 \\
colocated pairs & 75, across 84 stations \\
current baseline & \texttt{cls\_depth\_star\_reg}, 16 epochs, 10h13m, 4×H100 \\
current ubRMSE 0–10 & 0.0539 \\
current overfit & 21× — train 0.000099 vs val 0.002177 \\
\bottomrule
\end{tabularx}
\end{center}

The label count is abundant; the \emph{effective} sample size is not, and it differs
per parameter:

\begin{Verbatim}[fontsize=\scriptsize,frame=leftline,framerule=0.6pt,rulecolor=\color{rulegray},xleftmargin=6pt,framesep=6pt,]
   the shared dynamics  R(t), E(t), m, Q0       ~2.6 M station-days   abundant
   the map  description -> (lambda, alpha, beta)   ~993 locations     THE CONSTRAINT
   the within-tile contrast  f, pair loss           75 pairs          thin
\end{Verbatim}

One station per tile means anything about \emph{between-location} differences is a
993-sample problem, not a 2.6-million-sample one. This is the same wall §27b
hit, and it is the honest centre of the feasibility question.

\subsection*{0.2 Finding — trajectory output is a 365× win}
\addcontentsline{toc}{subsection}{0.2 Finding — trajectory output is a 365× win}

Today one sample is one station on one date, so a forward pass yields **one
label*\emph{. This design emits a }trajectory*: the water balance runs over the
365-day window the transformer already reads, so every day of that window comes
out at no extra cost.

\begin{Verbatim}[fontsize=\scriptsize,frame=leftline,framerule=0.6pt,rulecolor=\color{rulegray},xleftmargin=6pt,framesep=6pt,]
   labels per forward pass          1        ->   up to 365 x 3 = 1,095
   forward passes for full coverage 2.6 M    ->   7,164  (one per station-year)
\end{Verbatim}

Honest caveat: 7,164 samples at batch 128 is roughly **56 optimiser steps per
epoch**. Few steps, far richer gradients per step. The optimisation dynamics are
unknown and this is a genuine open question — mitigate with random window start
dates as augmentation, which multiplies the sample count without new data.

\subsection*{0.3 Finding — compute is a non-issue}
\addcontentsline{toc}{subsection}{0.3 Finding — compute is a non-issue}

The classical TOPMODEL result that makes this cheap: the catchment deficit
evolves on the \textbf{TWI histogram (\textasciitilde{}20 bins), not on the grid}. The saturated
fraction is an integral over the TWI distribution, so no per-cell loop is needed
for the far-field term at all.

\begin{Verbatim}[fontsize=\scriptsize,frame=leftline,framerule=0.6pt,rulecolor=\color{rulegray},xleftmargin=6pt,framesep=6pt,]
   TRAINING    365 steps x (20 TWI bins + 2 supervised cells)  ~ 8 k scalar ops/sample
   INFERENCE   365 steps x 4,900 cells                         ~ 1.8 M scalar ops
   AUTOGRAD    365 x B x cells x 2 stores x 4 bytes            < 1 MB at B=128
   OVERHEAD    ~5-10 % against a 50 M-param transformer over 1,035 tokens
\end{Verbatim}

The deleted U-Net decoder — four upsampling stages at 512/256/128/64 plus three
FiLM'd skips, \textasciitilde{}3.2 M activations per sample per depth — more than pays for the
loop. The only real engineering risk is that the loop is **365 sequential kernel
launches**; it is launch-bound, not memory-bound. Measure it in the smoke run
and chunk or \texttt{scan} it if it bites.

\subsection*{0.4 Finding — the 21× overfit is structurally capped}
\addcontentsline{toc}{subsection}{0.4 Finding — the 21× overfit is structurally capped}

This is the strongest feasibility argument, and it is structural rather than a
penalty.

\begin{Verbatim}[fontsize=\scriptsize,frame=leftline,framerule=0.6pt,rulecolor=\color{rulegray},xleftmargin=6pt,framesep=6pt,]
   TODAY   station identity  ->  decoder  ->  arbitrary value at any date
           measured overfit 21x;  §19.5 proved regularisation cannot touch it;
           the diagnosis was station memorisation

   HERE    station identity  ->  lambda, alpha, beta, p, delta  ->  water balance
                                 five bounded scalars
\end{Verbatim}

To memorise a 2,600-day series you would have to route it through five numbers
whose ranges are fixed by sigmoid and softplus. That is not discouraged — it is
impossible. The model's capacity to fit a station is bounded by the expressive
power of a two-store bucket, which is exactly the point.

\subsection*{0.5 Finding — a leak that must be closed, or this is worse than today}
\addcontentsline{toc}{subsection}{0.5 Finding — a leak that must be closed, or this is worse than today}

If \texttt{R(t)} and \texttt{E(t)} are read off \texttt{ctx[:, era5:]}, those day-tokens have already
attended to the station's own 196 spatial tokens. That supplies 365 × 2
station-specific, date-specific numbers — **reopening the memorisation channel
and defeating §0.4 entirely**.

\begin{Verbatim}[fontsize=\scriptsize,frame=leftline,framerule=0.6pt,rulecolor=\color{rulegray},xleftmargin=6pt,framesep=6pt,]
   REQUIRED CHANGE

     R(t), E(t) = MLP( [ raw ERA5 row(t) , g_static ] )        ~5 k params

     g_static = pooled soil / terrain / LULC summary ONLY.
                No per-date content. No station-identifying content. Ever.
\end{Verbatim}

This costs a small MLP rather than a second transformer, and it is not optional.
It was found by the feasibility pass, not by the design.

\subsection*{0.6 Finding — spin-up}
\addcontentsline{toc}{subsection}{0.6 Finding — spin-up}

\begin{Verbatim}[fontsize=\scriptsize,frame=leftline,framerule=0.6pt,rulecolor=\color{rulegray},xleftmargin=6pt,framesep=6pt,]
   SURFACE STORE   lambda ~ 0.9 forgets its initial condition in ~30 days
                   (0.9^30 = 0.04). A 365-day window is ample.          OK

   CATCHMENT S      recession scale m may be months, so 365 days of
                   spin-up may not be enough.                           RISK
                   Fix: S is ONE scalar. Run it on a 3-year ERA5 window
                   at negligible cost, or warm-start from a climatology
                   predicted from g. Cheap either way.
\end{Verbatim}

\subsection*{0.7 The decisive pre-GPU experiment}
\addcontentsline{toc}{subsection}{0.7 The decisive pre-GPU experiment}

The entire architecture rests on one untested assumption: **that a location's
water-balance parameters are predictable from what the landscape looks like.**
That is directly measurable on CPU, before any GPU is spent.

\begin{Verbatim}[fontsize=\scriptsize,frame=leftline,framerule=0.6pt,rulecolor=\color{rulegray},xleftmargin=6pt,framesep=6pt,]
  STEP 1   Calibrate a three-parameter bucket to each station's OWN daily
           series, forced by that station's own ERA5. 993 independent fits,
           Pool(64), sbatch, minutes.

           Yields:
             (a) fitted lambda-hat, alpha-hat, beta-hat per station
             (b) THE CEILING — how well can the best possible bucket fit at
                 all? If the best fit is worse than the current 0.0539 ubRMSE,
                 the physics is too crude and the approach caps below the
                 baseline. Hard stop, measured in an afternoon.
             (c) per-station equifinality, by profile likelihood

  STEP 2   Regress the fitted parameters on tokens + terrain + soil, reusing
           station_mean_probe.py's ladder: fit_block(), GroupKFold on
           location_group_id, RidgeCV as the deliberately weak learner,
           HistGradientBoostingRegressor as the nonlinearity control.

  STEP 3   Read the skill.
             > 30 %   the architecture has a foundation. Build it.
             ~ 10 %   marginal. This is where mean SM landed (§27b: 9.8 %).
             ~  0 %   the map does not exist. Do not build it.
\end{Verbatim}

This is the same shape as §29 Phase A and §32.8's sufficiency gate: a cheap
measurement that can kill the arm before the allocation is spent. It is worth
more than every other page of this document.

Note the target is \textbf{not} the one §27b tested. §27b regressed *mean soil
moisture*, which is largely climate and sensor calibration, and got 9.8 \%
globally and 0–4.2 \% within-network. Here the target is a \emph{drainage rate} — a
landscape property, which is precisely what §27a's positive controls proved the
tokens do encode — and its predictors include TWI, HAND, slope and texture as
direct measurements rather than embeddings. Different question, and it may have
a different answer. It may also not.

\subsection*{0.8 Data availability}
\addcontentsline{toc}{subsection}{0.8 Data availability}

\begin{Verbatim}[fontsize=\scriptsize,frame=leftline,framerule=0.6pt,rulecolor=\color{rulegray},xleftmargin=6pt,framesep=6pt,]
   ON DISK          labels; ERA5 365-day windows; soil (21,74,74) = 3 depths
                    x 7 properties; S1/S2 rasters at 10 m; TerraMind tokens
   IN PROGRESS      TWI / HAND — 353/353 regions fetched, build_twi_hand.py
                    running, MERIT gate (§32.6) not yet run
   NEVER MEASURED   per-station Sentinel-1 acquisition counts. Needed to size
                    the assimilation rung (V5) and nothing else.
   TO DOWNLOAD      nothing
\end{Verbatim}

\subsection*{0.9 What would make it infeasible}
\addcontentsline{toc}{subsection}{0.9 What would make it infeasible}

\begin{Verbatim}[fontsize=\scriptsize,frame=leftline,framerule=0.6pt,rulecolor=\color{rulegray},xleftmargin=6pt,framesep=6pt,]
   1.  §0.7 Step 1's bucket ceiling worse than 0.0539 ubRMSE     -> dead
   2.  §0.7 Step 3's skill ~ 0                                   -> dead
   3.  §32.8's sufficiency gate flat both wet and dry            -> V4 dies;
                                                                   V1-V3 may not
                                                                   beat baseline
   4.  56 steps/epoch too few to optimise                        -> fixable by
                                                                   window-start
                                                                   augmentation
\end{Verbatim}

Failures 1 and 2 are cheap and come first. Failure 3 prunes one input rather
than the architecture, because the per-cell parameters do not depend on terrain
alone.

\medskip\hrule\medskip

\section*{1. THE DIAGNOSIS}
\addcontentsline{toc}{section}{1. THE DIAGNOSIS}

The model reproduces \textbf{temporal} variation almost perfectly and \textbf{spatial}
variation barely at all.

\begin{Verbatim}[fontsize=\scriptsize,frame=leftline,framerule=0.6pt,rulecolor=\color{rulegray},xleftmargin=6pt,framesep=6pt,]
   within-station temporal SD          0.051   vs observed 0.051
   between-station SD                  15-19 % of observed
   r(pred level, obs level), CR200-18  -0.175          anti-correlated
   §23 norm-std                        0.093-0.524, HIGHEST at the flattest site
\end{Verbatim}

The cause is verified at \texttt{model.py:736-743}. \texttt{context} is the masked mean of
every valid non-CLS, non-spatial, non-pad token, applied by FiLM as a \textbf{uniform}
per-channel scale and shift. All time-series signal — ERA5's 365 days, the S2
and S1 histories, SIF, TWSA — reaches the decoder through **one spatially
constant vector**.

\begin{Verbatim}[fontsize=\scriptsize,frame=leftline,framerule=0.6pt,rulecolor=\color{rulegray},xleftmargin=6pt,framesep=6pt,]
   Two stations 500 m apart can be given different MEANS.
   They cannot be given different BEHAVIOUR.
\end{Verbatim}

That is exactly the observed signature: temporal skill perfect, between-station
skill absent. The map is not flat — it is wrong.

Three further measurements constrain any fix:

\begin{itemize}[leftmargin=1.2em,itemsep=1pt,topsep=3pt]
\item \textbf{§27a} — the tokens carry genuine within-tile spatial structure at 160 m,
and it tracks the landscape. Both positive controls passed: the LULC PCA→RGB
reproduces the land-cover map, the DEM PCA→RGB traces the drainage. Sink-
excluded PCA ratios are high-rank at every depth.
\item \textbf{§27b} — the same tokens do \emph{not} predict wetness. Own-token within-network
skill 0–4.2 \%; global 9.8 \% against a tabular stack's 23.4 \%. The UMAP
organises by Köppen class, not by soil moisture.
\item \textbf{§29.15} — the Landsat LST tile pattern is +0.967 spatially coherent across
all twelve months. It is a static landscape field, and *a static field cannot
track a dynamic variable*. That is why the pooled correlation was +0.167 and
the within-station correlation −0.077: the same fact seen twice.
\end{itemize}

The design principle follows directly. **Ask each source only for what it has
been measured to be good at.**

\begin{Verbatim}[fontsize=\scriptsize,frame=leftline,framerule=0.6pt,rulecolor=\color{rulegray},xleftmargin=6pt,framesep=6pt,]
   ERA5 weather      the dynamics, shared    temporal SD 0.051 = observed 0.051
   imagery/terrain   the PARAMETERS          tracks landscape (§27a) but not
                     not the wetness         wetness (§27b)
   soil/climate      the absolute level      23.4 % skill vs the network's ~18 %
   tables
   Sentinel-1        corrections to state    the only time-varying 10 m carrier
\end{Verbatim}

\medskip\hrule\medskip

\section*{2. THE FLOWCHART}
\addcontentsline{toc}{section}{2. THE FLOWCHART}

\begin{Verbatim}[fontsize=\scriptsize,frame=leftline,framerule=0.6pt,rulecolor=\color{rulegray},xleftmargin=6pt,framesep=6pt,]
+==============================================================================+
| OFFLINE  precomputed, on disk, never inside the training loop                |
+==============================================================================+

 TerraMind ViT-B (frozen)      §32 regional flow fields      soil rasters
 S2 S1 DEM LULC -> tokens      TWI_k  HAND_k  TWI_mean       OpenLandMap
 (N,196,768) fp16              at FULL catchment scale       -> theta_r, AWC,
        |                              |                        n_e, L
        +--------------------+---------+------------------------+
                             v
+==============================================================================+
| BLOCK A   DESCRIBE                              existing code, unchanged     |
+==============================================================================+

 seq = [ DEM x4 | LULC x4 | soil x4 | spatial x196 | S2hist | S1hist | ERA5 x365 ]
                                  ^
                                  |  terrain_stem, ADDED, zero-init
                                  |  conv stride-5 over (TWI, HAND, slope, curv)
                        +---------+----------+
                        |  6 x self-attention |
                        +---------+----------+
                                  |
                            ctx (B,1035,768)
                                  |
     +----------------------------+----------------------------+
     v                            v                            v
 ctx[:,12:208]              masked mean                  raw ERA5 row(t)
 196 cells                  g (B,768)                    + g_static
                                                         (NOT from ctx -- §0.5)

+==============================================================================+
| BLOCK B   PARAMETERISE                       four small heads, ~24 k params  |
+==============================================================================+

+-----------------+ +----------------+ +---------------+ +-------------------+
| B1  CELL        | | B2  FORCING    | | B3  READOUT   | | B4  CATCHMENT     |
|                 | |                | |               | |                   |
| gather x5       | | MLP([era5(t),  | | pedotransfer  | | MLP(g -> 2)       |
|  -> 70x70       | |      g_static])| | on soil       | |                   |
| + fine channels | |                | |               | |  m   mm           |
| MLP(768+C -> 5) | |  R(t)  mm/d    | |  theta_r      | |  Q0  mm/d         |
|                 | |  E(t)  mm/d    | |  AWC          | |  f   FIXED        |
|  lambda_k       | |                | |  n_e, L       | |                   |
|  alpha_k        | | SHARED over    | |  NOT LEARNED  | | ONE catchment     |
|  beta_k         | | all cells      | |               | | per tile          |
|  p_k            | |                | | + delta_d     | |                   |
|  K_k            | | (B,365,2)      | |   small, learn| | (B,2)             |
| (B,4900,5)      | |                | | (B,4900,3,2)  | |                   |
+--------+--------+ +-------+--------+ +-------+-------+ +---------+---------+
         +------------------+------------+--------+---------------+
                                         v
+==============================================================================+
| BLOCK C   SIMULATE                    365 days x 4,900 cells @ 32 m          |
+==============================================================================+

 for t = 1 .. 365:

 +---- FAR FIELD ---- the catchment as ONE state ----------------------------+
 |                                                                           |
 |   S(t)      catchment storage deficit, mm                                 |
 |   D_k(t)  = S(t) + (1/f)*( TWI_mean - TWI_k )        local deficit, mm    |
 |             |moves|        |---- fixed, from §32 ----|                    |
 |   z_k(t)  = D_k(t) / (1000*n_e)                      water table depth, m |
 |   u*_k(t) = exp( -z_k(t) / L_k )                     equilib. saturation  |
 |                                                                           |
 |   D_k <= 0  ==>  SATURATED: u1 = 1, incoming R becomes runoff             |
 +-----------------------------------+---------------------------------------+
                                     |  Phi_k = omega * max(0, u*_k - u_k)
                                     |          capillary rise, ONE-SIDED
                                     v
 +---- NEAR FIELD ---- per cell k -------------------------------------------+
 |                                                                           |
 |  u1(t+1) = clip[ lambda_k*u1  + alpha_k*R(t)  - beta_k*E(t)*u1            |
 |                                - p_k*u1       + Phi1_k                    |
 |                  ( + kappa_k*[DIV_W u1](k)   V6 only )   , 0, 1 ]         |
 |                                                                           |
 |  u2(t+1) = clip[ lambda2_k*u2 + p_k*u1 - r_k(t) + Phi2_k , 0, 1 ]         |
 +-----------------------------------+---------------------------------------+
                                     |
                       +-------------v--------------+
                       | on Sentinel-1 dates:       |
                       |  sigma_hat = h(u1, rough,  |
                       |                 veg)       |
                       |  u1 <- u1 + K_k*(sigma_obs |
                       |                - sigma_hat)|
                       +-------------+--------------+
                                     |
                dS/dt = Q0*exp(-S/m) - mean_k( r_k )
                        |baseflow out|   |recharge in|
                                     v
                       u1(k,t), u2(k,t)   all cells

+==============================================================================+
| BLOCK D   OBSERVE                                                            |
+==============================================================================+

   y_d(k,t) = theta_r(k,d) + AWC(k,d) * u^{l(d)}_k(t) + delta_d(k)

      l(0-10 cm)  = surface store u1
      l(10-30 cm) = deep store u2
      l(30-50 cm) = deep store u2

                      y-hat (B, 3, 70, 70)   daily, 32 m
                              |
        +---------------------+---------------------+
        v                     v                     v
   +----------+     +-----------------+    +----------------+
   | L_abs    |     | L_pair          |    | L_lst   (aux)  |
   | Huber at |     | Huber on        |    | Linear(S_k->1) |
   | station  |     | y_i - y_j       |    | vs Landsat     |
   | cell k*  |     | SAME DATE       |    | anomaly        |
   +----------+     +-----------------+    +----------------+
\end{Verbatim}

\medskip\hrule\medskip

\section*{3. BLOCK A — DESCRIBE}
\addcontentsline{toc}{section}{3. BLOCK A — DESCRIBE}

\subsection*{What it does}
\addcontentsline{toc}{subsection}{What it does}

Produces one 768-number description per location, and one per day. This is
\texttt{model.py} exactly as it is today: the same sequence, the same six
self-attention layers, the same weights. The only addition is \texttt{terrain\_stem},
a stride-5 convolution over the terrain rasters, added into the 196 spatial
tokens with \textbf{zero-initialised} weights.

\subsection*{Equations}
\addcontentsline{toc}{subsection}{Equations}

\begin{Verbatim}[fontsize=\scriptsize,frame=leftline,framerule=0.6pt,rulecolor=\color{rulegray},xleftmargin=6pt,framesep=6pt,]
   ctx  = Transformer( seq )                    seq in R^(B x 1035 x 768)

   S_k  = ctx[:, 12+k, :]                       description of cell k
   g    = masked_mean( ctx )                    description of the tile
\end{Verbatim}

Note what is \emph{not} here: the forcing is \textbf{not} read off \texttt{ctx}. See §0.5 — that
would reopen the memorisation channel §0.4 closes.

\subsection*{Parameters}
\addcontentsline{toc}{subsection}{Parameters}

\begin{center}
\small
\begin{tabularx}{\linewidth}{XXXX}
\toprule
\textbf{symbol} & \textbf{shape} & \textbf{source} & \textbf{meaning} \\
\midrule
\texttt{terrain\_stem} & \textasciitilde{}60 k & learned, zero-init & maps (TWI, HAND, slope, curvature) at 32 m into the 160 m token grid \\
\texttt{S\_k} & 768 & output & what cell k looks like, in context of the whole tile \\
\texttt{g} & 768 & output & what the tile looks like \\
\bottomrule
\end{tabularx}
\end{center}

\subsection*{Logic and foundation}
\addcontentsline{toc}{subsection}{Logic and foundation}

\textbf{Why terrain enters here, additively, zero-init.} Three consequences, each
load-bearing. Adding zero changes nothing on day 0, so the forward pass is
bit-identical to \texttt{best.pt} and §31.8's regression gate passes trivially. Token
\emph{k}'s terrain lands on token \emph{k}, whereas appending pooled terrain tokens the
\texttt{dem\_pyr} way would make it tile-level and able only to shift a mean. And \texttt{g}
receives terrain anyway through attention, because it is computed from the
encoder \emph{output}.

\textbf{Why this block is asked for a description and not a prediction.} §27a
established that the tokens carry real within-tile spatial heterogeneity at
160 m and that it tracks terrain and land cover — the two positive controls
both passed. §27b established that the same tokens do not track wetness. So
this block is asked for \emph{what kind of place this is}, which is measured to
work, and is not asked for \emph{how wet it is}, which is measured not to.

\medskip\hrule\medskip

\section*{4. BLOCK B — PARAMETERISE}
\addcontentsline{toc}{section}{4. BLOCK B — PARAMETERISE}

\subsection*{4.1 B1 — cell parameters}
\addcontentsline{toc}{subsection}{4.1 B1 — cell parameters}

\subsubsection*{What it does}
\addcontentsline{toc}{subsubsection}{What it does}

Converts each location's description into the properties that govern how it
holds water.

\subsubsection*{Equation}
\addcontentsline{toc}{subsubsection}{Equation}

\begin{Verbatim}[fontsize=\scriptsize,frame=leftline,framerule=0.6pt,rulecolor=\color{rulegray},xleftmargin=6pt,framesep=6pt,]
   [ lambda_k , alpha_k , beta_k , p_k , K_k ]  =  g_cell( [ S_{k//5} , fine_k ] )

   S_{k//5}   the 768-vector of the 160 m token containing cell k   (gather)
   fine_k     measured 32 m channels at cell k: TWI, HAND, slope, curvature,
              S1 sigma0 latest and median, soil texture, LULC
\end{Verbatim}

\subsubsection*{Parameters}
\addcontentsline{toc}{subsubsection}{Parameters}

\begin{center}
\footnotesize
\begin{tabularx}{\linewidth}{XXXXXX}
\toprule
\textbf{symbol} & \textbf{name} & \textbf{units} & \textbf{range} & \textbf{forced by} & \textbf{physical meaning} \\
\midrule
\texttt{lambda} & retention & — & (0,1) & sigmoid & fraction of water still present tomorrow under gravity drainage alone \\
\texttt{alpha} & inverse capacity & 1/mm & (0, 1/30) & softplus, capped & how much one millimetre of rain raises relative saturation. Small = deep or porous soil \\
\texttt{beta} & ET over capacity & 1/mm & (0, 1/20) & softplus, capped & how strongly atmospheric demand draws this cell down. Large = shallow soil, dense vegetation, exposed \\
\texttt{p} & percolation & 1/day & (0,1) & sigmoid & fraction of the surface store passed to the deeper store each day \\
\texttt{K} & radar gain & — & (0,1) & sigmoid & how much to trust a Sentinel-1 innovation here \\
\bottomrule
\end{tabularx}
\end{center}

\subsubsection*{Logic and foundation}
\addcontentsline{toc}{subsubsection}{Logic and foundation}

\textbf{Why three dynamic parameters and not the six an obvious design would use.}
Substituting \texttt{u = s/c} into the balance shows that capacity \texttt{c} and ET rate
\texttt{gamma} appear only as the combinations \texttt{R/c} and \texttt{gamma/c}:

\begin{Verbatim}[fontsize=\scriptsize,frame=leftline,framerule=0.6pt,rulecolor=\color{rulegray},xleftmargin=6pt,framesep=6pt,]
   s(t+1) = clip( lambda*s + R - E*gamma*s/c , 0, c )

   divide by c:

   u(t+1) = clip( lambda*u + R/c - E*(gamma/c)*u , 0, 1 )
                             |---|      |-------|
                            alpha         beta
\end{Verbatim}

\texttt{c} never appears alone. There is a one-dimensional family of \texttt{(c, gamma)}
giving identical predictions, so naming both is naming a parameter the data can
never see. Parameterising \texttt{alpha = 1/c} and \texttt{beta = gamma/c} removes an \textbf{exact}
degeneracy rather than hoping training resolves it. This reduction came out of
the equifinality critique and it makes the model smaller, not larger.

\textbf{Why every parameter is bounded by its activation.} The network cannot create
water (\texttt{lambda < 1}), cannot have negative capacity, and cannot evaporate what is
not there. It chooses inside a box that physics defines. Compare the U-Net, which
was free to paint anything and — per §23 — painted a field anti-correlated with
the landscape.

\textbf{Why a shared map rather than per-station parameters.} \texttt{g\_cell} is one MLP
applied at all 993 stations. That is the difference between this and classical
hydrological calibration, and it is the whole answer to equifinality of
predictions — see §10.2.

\subsection*{4.2 B2 — forcing}
\addcontentsline{toc}{subsection}{4.2 B2 — forcing}

\subsubsection*{What it does}
\addcontentsline{toc}{subsubsection}{What it does}

Converts each day's weather into two scalars: how much water enters the soil,
and how hard the atmosphere pulls.

\subsubsection*{Equation}
\addcontentsline{toc}{subsubsection}{Equation}

\begin{Verbatim}[fontsize=\scriptsize,frame=leftline,framerule=0.6pt,rulecolor=\color{rulegray},xleftmargin=6pt,framesep=6pt,]
   [ R(t) , E(t) ]  =  softplus( MLP( [ era5_row(t) , g_static ] ) )

   g_static = pooled soil / terrain / LULC summary.  Static. Tile-level.
              NO per-date content. NO station-identifying content.
\end{Verbatim}

\subsubsection*{Parameters}
\addcontentsline{toc}{subsubsection}{Parameters}

\begin{center}
\small
\begin{tabularx}{\linewidth}{XXXX}
\toprule
\textbf{symbol} & \textbf{name} & \textbf{units} & \textbf{physical meaning} \\
\midrule
\texttt{R(t)} & effective infiltration & mm/day & rainfall that actually enters the soil, after interception, runoff, and freezing \\
\texttt{E(t)} & atmospheric demand & mm/day & what would evaporate if water were unlimited — a potential evapotranspiration \\
\bottomrule
\end{tabularx}
\end{center}

\subsubsection*{Logic and foundation}
\addcontentsline{toc}{subsubsection}{Logic and foundation}

\textbf{Why not raw ERA5 precipitation.} ERA5 says 12 mm fell. How much enters
depends on whether it was one storm or twelve hours of drizzle, whether there is
a canopy intercepting it, whether the ground is already saturated, and whether
it is frozen. Conditioning on \texttt{g\_static} lets the head emit \emph{effective} forcing
rather than raw rainfall, which is what the water balance needs.

\textbf{Why it is shared across all cells.} ERA5-Land is 9 km; the tile is 2.24 km.
There is no within-tile weather information in the data, and inventing some
would be §23's failure repeated in a new place. This is the assumption the whole
design rests on: **the weather is shared, and the differences are in the
ground.**

\textbf{Why it must NOT be read off \texttt{ctx}.} This is §0.5, and it is the single most
important constraint in the document. Day-tokens inside \texttt{ctx} have attended to
the station's own spatial tokens, so reading forcing off them supplies
365 × 2 station-specific, date-specific numbers — a memorisation channel richer
than the one that produced the current 21× overfit. The whole benefit of §0.4
depends on this head being blind to per-date station content.

\subsection*{4.3 B3 — readout constants}
\addcontentsline{toc}{subsection}{4.3 B3 — readout constants}

\subsubsection*{What it does}
\addcontentsline{toc}{subsubsection}{What it does}

Converts relative saturation, which is dimensionless and lives in [0,1], into
volumetric soil moisture in m³/m³.

\subsubsection*{Equation}
\addcontentsline{toc}{subsubsection}{Equation}

\begin{Verbatim}[fontsize=\scriptsize,frame=leftline,framerule=0.6pt,rulecolor=\color{rulegray},xleftmargin=6pt,framesep=6pt,]
   y_d(k,t)  =  theta_r(k,d)  +  AWC(k,d) * u_k(t)  +  delta_d(k)
                |---------- pedotransfer ---------|    |-- learned --|
\end{Verbatim}

\subsubsection*{Parameters}
\addcontentsline{toc}{subsubsection}{Parameters}

\begin{center}
\footnotesize
\begin{tabularx}{\linewidth}{XXXXX}
\toprule
\textbf{symbol} & \textbf{name} & \textbf{units} & \textbf{source} & \textbf{physical meaning} \\
\midrule
\texttt{theta\_r} & residual water content & m³/m³ & pedotransfer on texture & what a probe reads when the plant-available store is empty — tightly bound water \\
\texttt{AWC} & available water capacity & m³/m³ & pedotransfer on texture and bulk density & the swing between wilting point and field capacity \\
\texttt{delta\_d} & correction & m³/m³ & learned, small, regularised toward zero & sensor calibration, install depth, and the point-versus-pixel support mismatch \\
\bottomrule
\end{tabularx}
\end{center}

\subsubsection*{Logic and foundation}
\addcontentsline{toc}{subsubsection}{Logic and foundation}

\textbf{Why \texttt{theta\_r} and \texttt{AWC} are read from tables, not learned.} Two independent
reasons. First, they are standard pedotransfer outputs computed from texture and
bulk density, and the rasters are already on disk as \texttt{soil (21,74,74)} — three
depths by seven properties. Second, and decisively for identifiability:
\texttt{y = A + B*u} means any rescaling of \texttt{u} can be absorbed by \texttt{B}. The only thing
pinning \texttt{u}'s scale is the clip at 1, which never engages at a site that never
saturates, so at a semi-arid station \texttt{B} and \texttt{alpha} are genuinely degenerate.
Fixing \texttt{B} from texture pins the scale unconditionally.

\textbf{Why the level goes to tables rather than to the network.} §27b's ladder
measured a tabular stack of soil, terrain, land cover, climate, lat/lon and SMAP
at 0.0752 → 0.0576, a skill of 23.4 \% on station-mean soil moisture. The trained
network's own RMS per-station bias is 0.0618 against the same 0.0752 null — about
18 \%. **The network is currently worse at predicting station level than
covariates already sitting on disk.** Delegating that job is a win by
measurement, not by argument.

\textbf{Why \texttt{delta\_d} exists and is kept small.} Part of every station's mean is
instrument, not landscape — calibration, installation depth, and the mismatch
between a point probe and a 32 m cell (§27b.6). No image contains that, and no
pedotransfer function does either. It is irreducible, so it gets an explicit,
regularised term rather than being allowed to contaminate every spatial
parameter in the model.

\subsection*{4.4 B4 — catchment}
\addcontentsline{toc}{subsection}{4.4 B4 — catchment}

\subsubsection*{What it does}
\addcontentsline{toc}{subsubsection}{What it does}

Carries the state of the whole upstream catchment as a single number, and
distributes it across the tile using the topographic index.

\subsubsection*{Equations}
\addcontentsline{toc}{subsubsection}{Equations}

\begin{Verbatim}[fontsize=\scriptsize,frame=leftline,framerule=0.6pt,rulecolor=\color{rulegray},xleftmargin=6pt,framesep=6pt,]
   S(t+1)  =  S(t)  +  Q0*exp( -S(t)/m )  -  mean_k( r_k(t) )      deficit, mm
                       |-- baseflow out --|   |-- recharge in --|

   D_k(t)  =  S(t)  +  (1/f)*( TWI_mean - TWI_k )                  local, mm
              |moves|         |------ fixed, from §32 ------|

   z_k(t)  =  D_k(t) / ( 1000 * n_e )                              depth, m

   u*_k(t) =  exp( -z_k(t) / L_k )                                 equilibrium

   D_k <= 0  ==>  saturated; u1 = 1 and incoming R becomes runoff
\end{Verbatim}

\subsubsection*{Parameters}
\addcontentsline{toc}{subsubsection}{Parameters}

\begin{center}
\footnotesize
\begin{tabularx}{\linewidth}{XXXXX}
\toprule
\textbf{symbol} & \textbf{name} & \textbf{units} & \textbf{source} & \textbf{physical meaning} \\
\midrule
\texttt{S} & catchment storage deficit & mm & \textbf{state}, evolved & how far the whole catchment is from saturation \\
\texttt{m} & recession scale & mm & learned from \texttt{g} & how fast baseflow declines as the catchment dries — the catchment's memory \\
\texttt{Q0} & baseflow at zero deficit & mm/day & learned from \texttt{g} & discharge when saturated \\
\texttt{f} & transmissivity decay & 1/mm & \textbf{fixed} at a literature value & how fast soil transmissivity falls with depth; sets the spatial spread of water-table depth \\
\texttt{n\_e} & drainable porosity & — & pedotransfer & converts a storage deficit into a depth \\
\texttt{L} & capillary fringe height & m & pedotransfer on texture & how far water rises above the water table. Sand 0.1–0.3, silt 1–2, clay 3–10 \\
\bottomrule
\end{tabularx}
\end{center}

\subsubsection*{Logic and foundation}
\addcontentsline{toc}{subsubsection}{Logic and foundation}

\textbf{The problem this solves.} A real catchment is tens of km²; the patch is
5 km². You cannot route the catchment — that would need container parameters
over 50 km, i.e. the transformer over 50 km. You cannot ignore it — the water
arrives regardless. §31.4 measured the consequence directly: **the lowest cell
lies on the tile boundary in 25 of 30 sampled tiles (83 \%)**, so flow leaves the
window and arrives from outside it.

\textbf{TOPMODEL's answer (Beven \& Kirkby 1979): never route it.} Carry one
catchment-average deficit and distribute it by the topographic index. All the
non-local information enters through \texttt{TWI\_k}, which §32 computes on the full
regional flow field and which is therefore correct at catchment scale by
construction. This is why §32's 353-region build is a dependency: the TWI it
produces is exactly the quantity this block needs, and a tile-local TWI would
not do.

\textbf{Why this makes terrain dynamic — the crux of the design.}

\begin{Verbatim}[fontsize=\scriptsize,frame=leftline,framerule=0.6pt,rulecolor=\color{rulegray},xleftmargin=6pt,framesep=6pt,]
   TWI AS A STATIC INPUT CHANNEL           TWI AS A DEFICIT DISTRIBUTION
   ----------------------------            -----------------------------
   feeds a network alongside slope,        D_k(t) = S(t) + (1/f)(TWI_mean - TWI_k)
   soil, land cover
                                           same fixed TWI, but the water table
   can only shift cell k's MEAN            it implies MOVES with the catchment

   the terrain effect is identical         the terrain effect is strong when wet
   in a drought and a flood                and vanishes when dry
\end{Verbatim}

\texttt{TWI\_k} is fixed; \texttt{S(t)} moves. So the water table under every cell rises and
falls together, but each cell sits at its own offset, set by how much of the
catchment drains through it. Same input, computed the same way by §32,
completely different functional role — and it is the role TOPMODEL derives it
for. §31.5 already states that *"TWI is a water-table-depth index, not a soil
moisture index"*; this uses it as one.

\textbf{Why Grayson's state-dependence falls out rather than being imposed.}

\begin{Verbatim}[fontsize=\scriptsize,frame=leftline,framerule=0.6pt,rulecolor=\color{rulegray},xleftmargin=6pt,framesep=6pt,]
   WET CATCHMENT, S small                 DRY CATCHMENT, S large

   many cells have D_k <= 0 and           no cell saturates; z_k is deep
   saturate; exp(-z/L) is large           everywhere; exp(-z/L) collapses

   -> the pattern organises by TWI        -> every cell is an independent 1-D
      "nonlocal control"                     column governed by texture and
                                             vegetation: "local control"
\end{Verbatim}

Grayson et al. (1997) and the Tarrawarra work (Western, Grayson \& Blöschl)
established that terrain control of surface moisture is \textbf{state-dependent}.
§31.5 argued this is what dictates the architecture. Here it is a \emph{consequence}
of one state crossing a threshold, spending no free parameters on it. An earlier
draft of this design used a learned sigmoid gate on tile wetness; saturation
does the same job with two fewer parameters and a derivation behind it.

\textbf{Why \texttt{f} is fixed rather than learned.} \texttt{m} is identified by the temporal data
— 993 daily series constrain a recession rate well. \texttt{f} is identified only by
the spatial data, which here is the 75 colocated pairs. Classical TOPMODEL's
documented equifinality comes from calibrating both against **one outlet
hydrograph**, a purely temporal signal, so \texttt{f} was never constrained. Here they
are constrained by different dimensions of the data — but the spatial dimension
is thinly sampled, so \texttt{f} is held at a literature value unless §32.8's
sufficiency gate demonstrates the pairs can identify it.

\medskip\hrule\medskip

\section*{5. BLOCK C — SIMULATE}
\addcontentsline{toc}{section}{5. BLOCK C — SIMULATE}

\subsection*{Equations}
\addcontentsline{toc}{subsection}{Equations}

\textbf{Surface store, 0–10 cm.}

\begin{Verbatim}[fontsize=\scriptsize,frame=leftline,framerule=0.6pt,rulecolor=\color{rulegray},xleftmargin=6pt,framesep=6pt,]
   u1(t+1) = clip[  lambda_k * u1(t)              gravity drainage
                  + alpha_k * R(t)                infiltration
                  - beta_k * E(t) * u1(t)         evapotranspiration
                  - p_k * u1(t)                   percolation to depth
                  + Phi1_k(t)                     capillary rise
                  ( + kappa_k * [DIV_W u1](k) )   hillslope routing, V6 only
                 ,  0 , 1 ]
\end{Verbatim}

\textbf{Deep store, 10–50 cm.}

\begin{Verbatim}[fontsize=\scriptsize,frame=leftline,framerule=0.6pt,rulecolor=\color{rulegray},xleftmargin=6pt,framesep=6pt,]
   u2(t+1) = clip[ lambda2_k * u2(t) + p_k * u1(t) - r_k(t) + Phi2_k(t) , 0 , 1 ]

   r_k(t) = recharge leaving the deep store to the water table
\end{Verbatim}

\textbf{Capillary exchange.}

\begin{Verbatim}[fontsize=\scriptsize,frame=leftline,framerule=0.6pt,rulecolor=\color{rulegray},xleftmargin=6pt,framesep=6pt,]
   Phi_k(t) = omega * max( 0 , u*_k(t) - u_k(t) )
\end{Verbatim}

\subsection*{Logic and foundation}
\addcontentsline{toc}{subsection}{Logic and foundation}

\textbf{Why the ET term is multiplied by \texttt{u}.} A dry soil resists evaporation. This
is the standard water-stress limitation, and without it a soil would evaporate
at potential rate down to zero, which no soil does.

\textbf{Why capillary rise is one-sided.} The water table can only \emph{supply} water
upward. Downward loss is already carried by \texttt{lambda} (gravity drainage) and \texttt{p}
(percolation), so a two-sided relaxation would double-count the loss and
over-dry any cell with a deep water table.

\textbf{Why two stores.} The deep store is a low-pass-filtered version of the surface
store, so the 10–30 and 30–50 cm bins come out naturally flatter and more
dominated by their level than 0–10 cm. That is what the observations do — §21.8
found 10–30 is the \emph{best} depth (ubRMSE 0.0500 vs 0.0539), and §24.12 found
identity dominating at depth. Both are consistent with a damped cascade.

\textbf{Why the clip matters beyond numerics.} \texttt{u} in [0,1] is what pins the scale of
the state. Together with \texttt{AWC} fixed from texture (§4.3), it is what stops the
readout and the dynamics trading off against each other.

\textbf{Why \texttt{DIV\_W} (V6) is optional and off by default.} TOPMODEL's \texttt{D\_k} already
handles lateral redistribution at every scale under a quasi-steady water table.
\texttt{DIV\_W} adds only the days-timescale transient in the unsaturated zone that
TOPMODEL assumes away. It also reintroduces the boundary problem this design was
built to avoid: a 70×70 patch has a wall at its edge, water arrives from
outside it, and perimeter cells shed into nothing so their \texttt{kappa} is
unidentifiable. Running V1–V4 first settles by measurement whether that
transient is worth the cost, and whether \texttt{DIV\_W} and \texttt{D\_k} are double-counting
the same within-patch accumulation.

\subsection*{Sentinel-1 assimilation}
\addcontentsline{toc}{subsection}{Sentinel-1 assimilation}

\begin{Verbatim}[fontsize=\scriptsize,frame=leftline,framerule=0.6pt,rulecolor=\color{rulegray},xleftmargin=6pt,framesep=6pt,]
   sigma_hat(k,t) = h( u1_k(t) , roughness_k , vegetation_k )     forward operator
   u1_k  <-  clip( u1_k + K_k * ( sigma_obs(k,t) - sigma_hat(k,t) ) , 0 , 1 )
                                  |------- innovation -------|
\end{Verbatim}

\textbf{Why the innovation and not the raw value.} Absolute backscatter is dominated
by surface roughness and vegetation structure, not by water. The *gap between
what the state predicted and what the radar saw* is mostly water. §30.3a reaches
for the same thing by feeding \texttt{(latest, long-term median)} and hoping the network
differences them; making it an explicit innovation does that structurally.

\textbf{Why staleness needs no embedding.} Between acquisitions the state coasts
forward on ERA5 alone. If there has been no pass for three weeks, the state has
not been corrected for three weeks and its error has grown on its own. The
staleness embedding at \texttt{model.py:505-514} becomes unnecessary on this path.

\medskip\hrule\medskip

\section*{6. BLOCK D — OBSERVE}
\addcontentsline{toc}{section}{6. BLOCK D — OBSERVE}

\subsection*{Equation}
\addcontentsline{toc}{subsection}{Equation}

\begin{Verbatim}[fontsize=\scriptsize,frame=leftline,framerule=0.6pt,rulecolor=\color{rulegray},xleftmargin=6pt,framesep=6pt,]
   y_d(k,t) = theta_r(k,d) + AWC(k,d) * u^{l(d)}_k(t) + delta_d(k)

   L  =  L_abs  +  w_p * L_pair  +  w_l * L_lst
\end{Verbatim}

\subsection*{The three losses}
\addcontentsline{toc}{subsection}{The three losses}

\textbf{\texttt{L\_abs}} — Huber at the station's own cell, per depth. This is today's loss
unchanged, with a per-sample cell index replacing the hardcoded \texttt{(112,112)} at
\texttt{model.py:779}.

\textbf{\texttt{L\_pair}} — Huber on \texttt{y\_i - y\_j} against the observed \texttt{y\_i - y\_j}, for pairs
of stations on the same date: 75 colocated pairs across 84 stations, plus every
same-ERA5-cell pair.

\begin{Verbatim}[fontsize=\scriptsize,frame=leftline,framerule=0.6pt,rulecolor=\color{rulegray},xleftmargin=6pt,framesep=6pt,]
   Today nothing in the objective ever compares two locations.

   Under an absolute per-pixel loss, the risk-minimising answer at an
   unsupervised location is the tile mean -- which is exactly what the model
   produces (between-station SD 15-19 % of observed).

   Differencing two stations on one date removes weather, season, climate,
   the ERA5 cell and the calibration epoch EXACTLY. It is §27b's Scale C
   construction applied to supervision instead of to analysis.
\end{Verbatim}

This is the only term that puts gradient directly on the quantity measured to
be broken. It is also the piece §31 lacks: §31 adds the \emph{capacity} for
locations to differ but leaves the objective indifferent to whether they do.

\textbf{\texttt{L\_lst}} — optional auxiliary. \texttt{Linear(S\_k -> 1)} predicting the Landsat
surface-temperature anomaly per token, tile mean removed for that date. Gives
Block A dense gradient — roughly 500 independent values per tile against 1,
since TIRS is 100 m (§31.6). Four caveats, all from §29 and §31.6: the pattern
is \textbf{static} (+0.967 month-to-month), so it disciplines the static branch only;
the 18.2 \% no-retrieval region contains the wettest station, so missingness is
correlated with the target; and if the target is a fixed map per station the
model may \textbf{memorise 993 maps} rather than learn structure — test LST-head
skill on validation against training stations in the first smoke run. Decayed
and ablatable; nothing downstream reads it.

\medskip\hrule\medskip

\section*{7. COMPLETE PARAMETER INVENTORY}
\addcontentsline{toc}{section}{7. COMPLETE PARAMETER INVENTORY}

\begin{center}
\footnotesize
\begin{tabularx}{\linewidth}{XXXXX}
\toprule
\textbf{symbol} & \textbf{per} & \textbf{units} & \textbf{source} & \textbf{identified by} \\
\midrule
\texttt{lambda} & cell & — & learned & recession shape, all 2.6 M station-days \\
\texttt{alpha} = 1/c & cell & 1/mm & learned & storm response amplitude \\
\texttt{beta} = gamma/c & cell & 1/mm & learned & dry-down rate under demand \\
\texttt{p} & cell & 1/day & learned & lag between depth bins \\
\texttt{K} & cell & — & learned & error drop on S1 acquisition dates \\
\texttt{R(t)}, \texttt{E(t)} & day & mm/day & learned & the whole dataset \\
\texttt{theta\_r} & cell, depth & m³/m³ & pedotransfer & not fitted \\
\texttt{AWC} & cell, depth & m³/m³ & pedotransfer & not fitted \\
\texttt{n\_e}, \texttt{L} & cell & —, m & pedotransfer & not fitted \\
\texttt{delta\_d} & station & m³/m³ & learned, regularised & station mean residual \\
\texttt{m} & catchment & mm & learned & \textbf{temporal} — 993 daily series \\
\texttt{Q0} & catchment & mm/day & learned & \textbf{temporal} — 993 daily series \\
\texttt{f} & catchment & 1/mm & \textbf{fixed} & \textbf{spatial} — only 75 pairs \\
\texttt{TWI}, \texttt{HAND} & cell & —, m & §32 & not fitted \\
\bottomrule
\end{tabularx}
\end{center}

\subsection*{New network parameters}
\addcontentsline{toc}{subsection}{New network parameters}

\begin{Verbatim}[fontsize=\scriptsize,frame=leftline,framerule=0.6pt,rulecolor=\color{rulegray},xleftmargin=6pt,framesep=6pt,]
   terrain_stem     ~60 k    zero-init
   B1 cell head     ~ 8 k
   B2 forcing head  ~ 5 k    (includes g_static projection)
   B3 delta head    ~ 2 k
   B4 catchment     ~ 2 k
   S1 operator h    ~10 k
                   -------
                   ~87 k     against a DELETED four-stage U-Net decoder
\end{Verbatim}

\subsection*{What is kept, changed, deleted}
\addcontentsline{toc}{subsection}{What is kept, changed, deleted}

\begin{Verbatim}[fontsize=\scriptsize,frame=leftline,framerule=0.6pt,rulecolor=\color{rulegray},xleftmargin=6pt,framesep=6pt,]
  KEPT UNCHANGED
    frozen TerraMind ViT; tokens on disk
    dataset.py sequence assembly (+ terrain added into the spatial tokens)
    the 6-layer fusion transformer and all its weights
    §32's terrain pipeline, exactly as being built

  CHANGED
    ctx[:, 12:208]    U-Net bottleneck        ->  five numbers per cell
    masked mean       FiLM on the decoder     ->  catchment params, delta
    ERA5 forcing      averaged into `context` ->  a separate small MLP that is
                                                  BLIND to per-date station
                                                  content (§0.5)
    the sample        one station-date        ->  one station x 365-day window,
                                                  emitting a trajectory

  DELETED
    UNetDecoder -- four upsampling stages at 512/256/128/64 + three FiLM'd skips
    the staleness embedding on this path
    (an earlier draft's learned wet/dry gate -- saturation replaces it)
\end{Verbatim}

\medskip\hrule\medskip

\section*{8. BUILD LADDER}
\addcontentsline{toc}{section}{8. BUILD LADDER}

Each rung is a flag with exactly one check. **A rung that fails its check is
deleted, not tuned.**

\begin{center}
\small
\begin{tabularx}{\linewidth}{XXXX}
\toprule
\textbf{} & \textbf{adds} & \textbf{params} & \textbf{earns its place if} \\
\midrule
\textbf{V0} & §0.7's CPU experiment & 0 & bucket ceiling beats 0.0539 \textbf{and} parameters are predictable from the landscape. \textbf{Gate: everything below is conditional on this.} \\
\textbf{V1} & per-cell bucket \texttt{lambda, alpha, beta}, one store & \textasciitilde{}8 k & across-cell SD of \texttt{y-hat} rises off zero; \texttt{lambda} correlates positively with TWI \\
\textbf{V2} & \texttt{theta\_r}, \texttt{AWC} from pedotransfer + \texttt{delta} & \textasciitilde{}2 k & r(pred level, obs level) goes from −0.175 to positive \\
\textbf{V3} & second store, \texttt{p} & \textasciitilde{}1 k & 10–30 and 30–50 improve without hurting 0–10 \\
\textbf{V4} & TOPMODEL \texttt{S}, \texttt{D\_k}, capillary coupling & \textasciitilde{}3 k & r(pred, TWI) \textbf{strong when wet, near zero when dry} \\
\textbf{V5} & Sentinel-1 assimilation & \textasciitilde{}10 k & error drops on days near an acquisition and not otherwise \\
\textbf{V6} & near-field \texttt{DIV\_W} routing & \textasciitilde{}1 k & anything left after V4 — probably nothing \\
\bottomrule
\end{tabularx}
\end{center}

\begin{Verbatim}[fontsize=\scriptsize,frame=leftline,framerule=0.6pt,rulecolor=\color{rulegray},xleftmargin=6pt,framesep=6pt,]
   V1 + V2 is a complete, coherent architecture on its own, and it is a
   SMALLER change to model.py than deleting the U-Net already requires.

   V4 is the intellectually interesting rung -- it is what makes terrain
   dynamic and connects to §31.5's theory -- but it should be EARNED by
   V1 working, not assumed.
\end{Verbatim}

\medskip\hrule\medskip

\section*{9. GATES}
\addcontentsline{toc}{section}{9. GATES}

\subsection*{Before any GPU}
\addcontentsline{toc}{subsection}{Before any GPU}

\begin{Verbatim}[fontsize=\scriptsize,frame=leftline,framerule=0.6pt,rulecolor=\color{rulegray},xleftmargin=6pt,framesep=6pt,]
  0.  §0.7's CPU experiment. Bucket ceiling AND parameter predictability.
      This gate can kill the whole design in an afternoon and must run first.

  1.  All new heads zero-init  =>  forward bit-identical to best.pt.
      Reuse the §26 provenance gate UNCHANGED: >= 99.5 % of rows
      bit-identical AND max |delta| <= 1 bf16 ULP (0.000977).
      DO NOT re-tighten it.

  2.  Mass balance: sum(in) - sum(out) - delta(storage) = 0 to float precision.

  3.  Synthetic catchment with analytic TWI: the S recession must match the
      analytic exponential; a plane, a cone and a V-valley must give the TWI
      values §32.5 Tier 1 already specifies (including §32.9.2's correction
      that a cone gives a = r/2 exactly, not one cell).

  4.  §32.8's sufficiency gate: regress observed delta-SM on delta-TWI and
      delta-HAND across the 75 pairs under station fixed effects, SPLIT
      WET/DRY. This decides both whether V4 is worth building and whether f
      is identifiable.
\end{Verbatim}

\subsection*{After training}
\addcontentsline{toc}{subsection}{After training}

\begin{Verbatim}[fontsize=\scriptsize,frame=leftline,framerule=0.6pt,rulecolor=\color{rulegray},xleftmargin=6pt,framesep=6pt,]
  5.  S(t) must be seasonal -- low in spring, high in late summer.
      Flat => the catchment state is inert and V4 is dead weight.

  6.  r(prediction, TWI) split by S terciles: STRONG when wet, NEAR ZERO when
      dry. A flat profile is a static offset dressed as terrain and MUST be
      reported as a failure, not as increased variance. This is the same
      statistic as gate 4, so the gate and the architecture test ONE claim.

  7.  Across-cell SD of y-hat, logged every epoch. Near-zero => collapsed to
      the tile mean -- the failure that looks like success until you plot a map.

  8.  IDENTIFIABILITY. Train 3 seeds; correlate the parameter MAPS, not the
      predictions.
          r > 0.9  => identified. Interpret it, plot it against TWI, publish it.
          r < 0.5  => NOT identified. Predictions may stand; every interpretive
                      claim about that parameter is dropped.
      Plus a flat-direction check: perturb each parameter +/-20 % and measure
      the loss response. A flat direction names the combination to reparameterise
      away, as `c` and `gamma` were in §4.1.

  9.  m, Q0 in physically sane ranges. m = 500 mm means it is fitting noise,
      whatever the RMSE says.

 10.  Skill targets, unchanged from §30.9:

          own-centre ubRMSE (0-10)                0.0301  ->  <=
          off-centre ubRMSE                       0.0345  ->  materially closer
          between-station spread, % of observed   15-19 % ->  > 35 %
          r(pred level, obs level), CR200-18      -0.175  ->  > 0

 11.  PHYSICALITY, not just variance. norm-std must DROP at SCAN Crossroads --
      the flat control currently receiving the most painted variation (0.52).
      A run that increases spread while r(anomaly, terrain) stays near zero has
      made the map more variable, not more correct, and must not be reported as
      success. §23 is explicit on this.

 12.  Report 14x14 and 70x70 SEPARATELY. If 70x70 helps only via the static
      branch, that is a resolution-limited offset, not resolved dynamics.
\end{Verbatim}

\medskip\hrule\medskip

\section*{10. KNOWN WEAKNESSES, STATED BEFORE BUILDING}
\addcontentsline{toc}{section}{10. KNOWN WEAKNESSES, STATED BEFORE BUILDING}

\subsection*{10.1 §27b's ceiling may be real}
\addcontentsline{toc}{subsection}{10.1 §27b's ceiling may be real}

If the imagery genuinely carries no information about how a place holds water,
B1 outputs near-constants and every cell in the tile predicts one curve. **No
architecture manufactures information.** What this design does is fail
\emph{legibly} — a flat \texttt{lambda} map says so directly, where a decoder's wrong map
looked plausible for two years.

§0.7 is designed to find this out before the allocation is spent.

\subsection*{10.2 Equifinality of parameters is NOT cured}
\addcontentsline{toc}{subsection}{10.2 Equifinality of parameters is NOT cured}

This distinction matters and should not be blurred:

\begin{Verbatim}[fontsize=\scriptsize,frame=leftline,framerule=0.6pt,rulecolor=\color{rulegray},xleftmargin=6pt,framesep=6pt,]
   EQUIFINALITY OF PREDICTIONS      largely cured by regionalisation. Classical
                                    calibration fits a free parameter vector to
                                    ONE catchment's hydrograph, which is
                                    underdetermined. Here theta_k = MLP(landscape)
                                    is ONE map that must work at 993 stations, so
                                    a compensating error must be a CONSISTENT
                                    function of slope, TWI, texture and land cover
                                    or it breaks somewhere else.

   EQUIFINALITY OF PARAMETERS       NOT cured. The MLP can settle anywhere on a
                                    compensating manifold, and which point it
                                    picks is arbitrary.

   WHAT IS THEREFORE AT RISK        gates 6 and 8 -- the interpretive story --
                                    NOT the soil moisture predictions.
\end{Verbatim}

The failure mode is: **the model works, and the story about why is
unfalsifiable.** Gate 8 decides whether the story may be told. That is a real
scientific cost and it should not be papered over.

\subsection*{10.3 The near-field boundary problem is real, and V6 reintroduces it}
\addcontentsline{toc}{subsection}{10.3 The near-field boundary problem is real, and V6 reintroduces it}

A 2.24 km patch does not contain a catchment. §31.4 measured the lowest cell on
the tile boundary in 25 of 30 tiles. V1–V5 avoid this entirely by never routing
water across the patch — TOPMODEL carries the catchment as one number instead.
V6 puts it back, and with it: a wall at the patch edge, missing external inflow,
and perimeter cells whose \texttt{kappa} is unidentifiable because they shed into
nothing. That is why V6 is last and off by default.

\subsection*{10.4 TOPMODEL's quasi-steady assumption}
\addcontentsline{toc}{subsection}{10.4 TOPMODEL's quasi-steady assumption}

TOPMODEL assumes the water table equilibrates instantly with the catchment
deficit. In flashy catchments that is wrong. V6 is the patch, and V6 has §10.3's
problem. There is no clean answer to this within the design; it is a known
limit of the approach, and gate 5 will show it as a poorly-behaved \texttt{S(t)}.

\subsection*{10.5 Getting the physics wrong is a failure mode a black box does not have}
\addcontentsline{toc}{subsection}{10.5 Getting the physics wrong is a failure mode a black box does not have}

A misspecified capillary term is a bug that training cannot route around. §31's
per-location transformer cannot be wrong about physics because it claims none.
That is the trade being made here, deliberately: **this model states its
mechanism, so the mechanism can be wrong in a detectable way.**

\subsection*{10.6 The optimisation regime is untested}
\addcontentsline{toc}{subsection}{10.6 The optimisation regime is untested}

\textasciitilde{}56 steps per epoch at batch 128 with 1,095 labels per sample is a different
regime from the current one. Rich gradients, few steps. Nothing in this
codebase's history speaks to it. Mitigation is random window-start
augmentation, which is free, but the risk is real and belongs in the smoke run.

\subsection*{10.7 It will not produce structure below 32 m}
\addcontentsline{toc}{subsection}{10.7 It will not produce structure below 32 m}

And must not be presented as doing so. §23's verdict — *"do not present this
model as producing 10 m soil-moisture maps"* — applies to 32 m here just as
firmly. Blockiness at 5×5 cells within a token is a \textbf{diagnostic}, not a
defect: it is the map saying the fine channels are contributing nothing.
Bilinear smoothing hides exactly that failure behind a plausible ramp, which is
how the current map looked reasonable while being anti-correlated with the
landscape. Do not smooth it. Report it.

\subsection*{10.8 \texttt{f} rests on 75 pairs}
\addcontentsline{toc}{subsection}{10.8 \texttt{f} rests on 75 pairs}

Held fixed at a literature value unless gate 4 says otherwise.

\medskip\hrule\medskip

\section*{APPENDIX A — the loop, with numbers}
\addcontentsline{toc}{section}{APPENDIX A — the loop, with numbers}

Two cells. Same weather. Illustrative parameters, TOPMODEL coupling omitted for
clarity; the lateral term shown here is V6's \texttt{DIV\_W}, used because it makes the
mechanism visible in three lines of arithmetic.

\begin{Verbatim}[fontsize=\scriptsize,frame=leftline,framerule=0.6pt,rulecolor=\color{rulegray},xleftmargin=6pt,framesep=6pt,]
                       capacity   retention   ET     shedding
   ---------------------------------------------------------------
   HOLLOW               100 mm      0.95      0.8      0.02      holds
   RIDGE                 60 mm      0.80      1.0      0.15      sheds

   The hollow has three ridge-like cells flowing into it. Gate at 0.6.
\end{Verbatim}

\textbf{Day 1 — 20 mm rain, 3 mm demand.} Start: hollow 40 mm, ridge 20 mm.

\begin{Verbatim}[fontsize=\scriptsize,frame=leftline,framerule=0.6pt,rulecolor=\color{rulegray},xleftmargin=6pt,framesep=6pt,]
   RIDGE      s = 20,  fullness 0.333        HOLLOW     s = 40,  fullness 0.400
     rain          + 20.00                     rain          + 20.00
     evaporation   -  1.00                     evaporation   -  0.96
     drains down   -  4.00                     drains down   -  2.00
     sheds sideways-  1.80                     receives      +  4.92
                   --------                                  --------
     s(day 2)      =  33.20  fullness 0.553    s(day 2)      =  61.96  f 0.620
\end{Verbatim}

\textbf{Day 2 — no rain, 5 mm demand.}

\begin{Verbatim}[fontsize=\scriptsize,frame=leftline,framerule=0.6pt,rulecolor=\color{rulegray},xleftmargin=6pt,framesep=6pt,]
   RIDGE      s = 33.20                       HOLLOW     s = 61.96
     evaporation   -  2.77                      evaporation   -  2.48
     drains down   -  6.64                      drains down   -  3.10
     sheds sideways-  2.99                      receives      +  8.22
                   --------                                   --------
     s(day 3)      =  20.80  fullness 0.347     s(day 3)      =  64.60  f 0.646
                            FELL 37 %                                 ROSE
\end{Verbatim}

The hollow keeps rising the day after rain while the ridge collapses, because
water is still arriving from upslope. That lag is real hydrology and it is
exactly what one shared weather vector structurally cannot produce.

\textbf{Through Block D:}

\begin{Verbatim}[fontsize=\scriptsize,frame=leftline,framerule=0.6pt,rulecolor=\color{rulegray},xleftmargin=6pt,framesep=6pt,]
                LEVEL    RANGE    fullness      soil moisture
   ----------------------------------------------------------
   HOLLOW       0.15  +  0.20  x   0.646    =      0.279
   RIDGE        0.12  +  0.16  x   0.347    =      0.176
                                                ---------
                                     spread  =      0.103
\end{Verbatim}

Against an observed within-tile spread of \textasciitilde{}0.108 at CR200-18, where the current
model produces \textasciitilde{}0.021.

\medskip\hrule\medskip

\section*{APPENDIX B — comparison with §31}
\addcontentsline{toc}{section}{APPENDIX B — comparison with §31}

\begin{center}
\small
\begin{tabularx}{\linewidth}{XXXX}
\toprule
\textbf{} & \textbf{today} & \textbf{§31 per-location processor} & \textbf{this} \\
\midrule
frozen ViT & yes & yes & yes \\
fusion transformer & 1× & 1× & 1× \\
per-location model & \textbf{none} & transformer, run 196× at inference & 3–5 scalars + a recurrence \\
forcing & inside the fusion sequence, averaged into one vector & \texttt{W}, cross-attended per location & a small MLP, blind to station identity \\
decoder & U-Net, 4 stages & deleted & deleted \\
new parameters & — & \textasciitilde{}2 M & \textasciitilde{}87 k \\
labels per forward pass & 1 & 1 & up to 1,095 \\
terrain's role & — & static input channel & distribution function of a moving deficit \\
what stops memorisation & nothing (21× overfit) & nothing structural & 5 bounded scalars \\
interpretable & no & no & yes — and therefore falsifiable \\
can be wrong about physics & no & no & \textbf{yes} \\
\bottomrule
\end{tabularx}
\end{center}

§31 and this design agree on the diagnosis and on deleting the U-Net. They
disagree on what the per-location object should be: a 768-dimensional
transformer output, or a handful of bounded physical scalars. Under one label
per tile across 993 stations, the second is a well-posed estimation problem and
the first is not — that is the whole argument.

\medskip\hrule\medskip

\section*{REFERENCES}
\addcontentsline{toc}{section}{REFERENCES}

\begin{itemize}[leftmargin=1.2em,itemsep=1pt,topsep=3pt]
\item Beven, K. J. \& Kirkby, M. J. (1979). A physically based, variable contributing
area model of basin hydrology. \emph{Hydrological Sciences Bulletin} 24(1), 43–69.
\item Beven, K. J. (2006). A manifesto for the equifinality thesis. *Journal of
Hydrology* 320, 18–36.
\item Grayson, R. B., Western, A. W., Chiew, F. H. S. \& Blöschl, G. (1997). Preferred
states in spatial soil moisture patterns: local and nonlocal controls.
\end{itemize}
  \emph{Water Resources Research} 33(12), 2897–2908.
\begin{itemize}[leftmargin=1.2em,itemsep=1pt,topsep=3pt]
\item Nobre, A. D. et al. (2011). Height Above the Nearest Drainage — a
hydrologically relevant new terrain model. \emph{Journal of Hydrology} 404, 13–29.
\item Vachaud, G., Passerat de Silans, A., Balabanis, P. \& Vauclin, M. (1985).
Temporal stability of spatially measured soil water probability density
function. \emph{Soil Science Society of America Journal} 49(4), 822–828.
\item Western, A. W., Grayson, R. B. \& Blöschl, G. (2002). Scaling of soil moisture:
a hydrologic perspective. \emph{Annual Review of Earth and Planetary Sciences} 30,
149–180.
\end{itemize}

Internal: §23, §24.12, §26.11, §27a, §27b, §29.13, §29.15, §30.3a, §31.4,
§31.5, §31.6, §31.8, §32.3, §32.6, §32.8, §32.9 of \texttt{text/training\_runbook.md}.

\end{document}
